% Ad Atticum 8.11 --- headnote
% ad-atticum-08-11, 1 March 49 BC, Formian villa

\textit{Cicero to Atticus, written from the Formian
villa on the Kalends of March 49 BC (the manuscript
dateline: \textit{Scr.\ in Formiano in K.\ Mart.\ a.\
705 (49)}). The most considered of the letters of this
fortnight, and the one that brings the political crisis
into contact with Cicero's own philosophical work. With
Caesar racing south through Apulia and Pompey poised at
Brundisium to cross to Greece, Cicero turns back to
\textit{De Re Publica} and the portrait of the
\textit{moderator rei publicae} drawn there --- the
``ruler of the commonwealth'' whose end is the
prosperous and honourable life of the citizens --- and
measures the actual leadership of both men against it.
Neither, he concludes, is governed by that goal; both
want to reign. \textit{Dominatio quaesita ab utroque
est}.}

\textit{Section 1 opens with the studied composure: he
is agitated but not so agitated as Atticus thinks, and
he is spending his time considering the figure of the
ideal statesman ``drawn out with sufficient care'' in
his own books. The quotation that follows is Scipio's
in the fifth book of \textit{De Re Publica} --- helmsman
to a fair course, physician to health, general to
victory, ruler of the commonwealth to the prosperous
life of the citizens --- and the ``bringer-to-completion''
of that work is the figure Cicero now no longer sees in
Roman public life. Section 2 brings Pompey's conduct
into focus: not retreat under necessity but design from
the first --- to set every land and sea in motion, to
bring in barbarian kings and savage tribes, to assemble
the largest of armies. This is the \textit{genus
Sullani regni} (``Sullan kingship''), long the aim of
many of those around him. Section 3 is the prophecy ---
not as a Cassandra-figure (``that one whom no one
believed''), but by reasoned conjecture --- an Iliad of
evils \textit{[Greek: Ilias]} hangs over Italy. Section
4 makes that vision concrete: Italy trodden underfoot
next summer by slave conscripts gathered from every
type. Sections 5--7 close on the daily business ---
Caesar's letter, the younger Balbus's mission to
Lentulus, Pompey's two careless letters which Cicero
has answered with care and forwards in copy, and a
request that Atticus send him the book of Demetrius of
Magnesia ``On Concord.'' The closing line --- \textit{vides
quam causam mediter} --- shows where the writing has
been tending all along.}
